\section*{Chapitre \ref{ch:fm:case-studies-ii} --- Études de cas II}

\begin{solution}{exo:fm:case-studies-ii:1}
$97.9/(97.9+266.7)=26.9~\%$~; $400/266.7=\$1.50$ par action, alors que la valeur comptable par action après la perte restait de 10.80~\$.
\end{solution}

\begin{solution}{exo:fm:case-studies-ii:2}
Voir la \cref{def:fm:case-studies-ii:dynamic}. La marge initiale statique est fixée en dollars à l'ouverture~: à mesure que la valeur de la position augmente, les mêmes dollars
ne représentent plus qu'une part moindre de celle-ci, de sorte que l'effet de levier du client auprès du courtier croît exactement avec la concentration.
\end{solution}

\begin{solution}{exo:fm:case-studies-ii:3}
$10\,000\times(80\,000-48\,078)=319$ millions de dollars~; la trésorerie s'épuise à $48\,078+300\,000\,000/10\,000=\$78\,078$.
\end{solution}

\begin{solution}{exo:fm:case-studies-ii:4}
Voir la démonstration de la \cref{prop:fm:case-studies-ii:zero}. Un standstill assorti d'une vente conjointe plus lente laisse la baisse totale au plus inchangée et la répartit par
accord~; sans celui-ci, chaque créancier a intérêt à vendre le premier, ce qui reporte sa part de la baisse sur les autres.
\end{solution}

\begin{solution}{exo:fm:case-studies-ii:5}
$0.531\times0.5/7=3.8~\%$ pour le premier, $0.531\times5.5/7=41.8~\%$ pour le sixième, $0.531\times6.5/7=49.3~\%$ pour le dernier~; un unique courtier vendant la
position entière subit une baisse moyenne de $0.531/2=26.6~\%$.
\end{solution}

\begin{solution}{exo:fm:case-studies-ii:6}
Une limite sur les soldes détenus auprès d'une même plateforme, un balayage quotidien de l'excédent vers une banque ou un \omterm{def:fm:the-asset-manager-and-the-fund:admin}{dépositaire}, et l'exigence de preuves de ségrégation. Les
conditions de service et états financiers de la bourse, ainsi que toute attestation de réserves face aux passifs, auraient été les documents à réclamer~;
les retards de retrait constituent le signal tardif.
\end{solution}

\begin{solution}{exo:fm:case-studies-ii:7}
Les trois premières places ne perdent rien~; la quatrième à la septième perdent 1.12, 2.41, 3.70 et 4.99 milliards de dollars. La dernière perd encore 4.99 milliards de dollars.
\end{solution}

\begin{solution}{exo:fm:case-studies-ii:8}
Dans une course, la perte d'un courtier dépend de son rang et de ce que le client détient ailleurs~: dans le modèle calibré, le dernier vendeur a besoin de 49.3~\%, et une
marge de 20~\% perd encore 4.99 milliards de dollars. Une marge doit croître avec la concentration, rapportée au volume du marché, et avec ce que le courtier ne voit pas.
\end{solution}

\begin{solution}{pb:fm:case-studies-ii:1}
\begin{enumerate}
\item Voir la \cref{def:fm:case-studies-ii:dynamic}.
\item Voir la \cref{def:fm:case-studies-ii:race}.
\item Voir la \cref{prop:fm:case-studies-ii:zero}.
\item Ensemble, ils supportent une baisse fixe, plus faible s'ils vendent lentement~; isolément, chacun peut reporter sa part sur les autres en vendant le premier.
\item Voir l'\cref{dat:fm:case-studies-ii:record}~: plus de 460 millions de dollars perdus en 45 minutes le 1\up{er} août 2012~; 400 millions de dollars d'actions privilégiées convertibles
en environ 73~\% des actions vendues le 6 août, sans vote des actionnaires~; la fusion avec GETCO le 1\up{er} juillet 2013 à 3.75~\$ par action ou un tiers
d'une nouvelle action.
\item 120 milliards de dollars d'exposition brute et 9--10 milliards de dollars de \omterm{def:fm:the-proprietary-market-making-firm:ownfunds}{fonds propres}, et une demande de standstill~; ils ont refusé, et vendu en partie par transactions de bloc le 26
mars, trois dans le cadre d'une liquidation gérée jusqu'en avril.
\item Des actifs représentant 6.6~\% des dettes envers les détenteurs de jetons liquides et 22.6~\% de l'ensemble des dettes au moment de la requête~; les clients ont retiré un montant net de 7.0 milliards de dollars du 1\up{er} au 11
novembre, dont 3.3 milliards de dollars le 7, et les retraits ont été suspendus le 8.
\item L'établissement a suspendu la négociation et annulé toutes les transactions depuis minuit le 8 mars~; les tribunaux ont jugé cela légitime, et la FCA lui a infligé une amende de
9\,245\,900~livres pour manquement à l'obligation de négociation ordonnée.
\item Le \Cref{tab:fm:case-studies-ii:rescue}~: anciens actionnaires passant de 100~\% à 26.9~\%, valeur comptable par action de 15.29~\$ à 10.80~\$ puis à 4.00~\$.
\item Sept courtiers de 17 milliards de dollars chacun, impact permanent linéaire, le coefficient 0.531 étant choisi de sorte que le sixième, avec 9.4~\%, perde 5.5 milliards de dollars.
\item 0, 0.34, 1.63, 2.92, 4.21, 5.50 et 6.79 milliards de dollars.
\item 0, 0, 0, 1.12, 2.41, 3.70 et 4.99 milliards de dollars~; 3.8~\%, 11.4~\%, 19.0~\%, 26.6~\%, 34.2~\%, 41.8~\% et 49.3~\%.
\item Expositions inégales, baisse du cours avant défaut, positions courtes et couvertures, transactions de bloc dont l'effet sur le cours affiché n'est pas linéaire, et
le recouvrement sur les actifs du client.
\item 319 millions de dollars à 80\,000~\$, 533 millions de dollars au plus haut~; une trésorerie de 300 millions de dollars s'épuise à 78\,078~\$.
\item Des marges dynamiques assorties de majorations pour concentration et pour liquidité, recalculées chaque jour, et le droit de les relever rapidement.
\item Exiger contractuellement du client qu'il déclare ses positions auprès de tous les courtiers~; suivre les dépôts publics~; supposer le pire lorsqu'il refuse de répondre.
\item Une limite par plateforme et un balayage quotidien de l'excédent, avec des limites plus strictes là où la ségrégation ne peut être vérifiée.
\item Détenir du capital pour une matinée comparable à celle du 1\up{er} août 2012, sous une limite stricte d'exposition agrégée, et organiser un financement confirmé pouvant être mobilisé en une journée.
\item Avec une marge de 9.4~\%, 0 pour le premier vendeur, jusqu'à 6.79 milliards de dollars pour le dernier (4.99 milliards de dollars à 20~\%)~; les anciens actionnaires ont conservé 26.9~\% de la
société de 2012, 391 millions de dollars de \omterm{def:fm:the-proprietary-market-making-firm:ownfunds}{fonds propres} comptables au lieu de 1\,057 millions de dollars.
\item Majorer dynamiquement la marge des clients concentrés, pour la sortie de la position entière plutôt que pour la part d'un seul courtier, et convenir à l'avance de la conduite à tenir si le
client fait défaut. Plafonner les soldes propres de la société sur toute plateforme, et détenir du capital et un financement confirmé couvrant une perte équivalant à une journée de sa plus forte exposition.
\end{enumerate}
\end{solution}

\begin{solution}{iq:fm:case-studies-ii:1}
La marge statique est fixée à l'ouverture en pourcentage du notionnel~; la marge dynamique est recalculée à partir de la volatilité courante du portefeuille,
de la concentration, de la liquidité et du biais.
\iqlookfor{érosion de la marge statique à mesure que les cours évoluent.}
\end{solution}

\begin{solution}{iq:fm:case-studies-ii:2}
Des limites pré-trade, une limite d'exposition agrégée assortie d'un kill switch, et des contrôles de déploiement~; le coût pour les propriétaires dépend du capital et du financement
détenus face à une telle perte, et de savoir si un sauvetage doit être organisé en quelques jours.
\iqlookfor{contrôles, puis capital.}
\end{solution}

\begin{solution}{iq:fm:case-studies-ii:3}
Peser la course~: vendre en premier reporte la perte sur les autres, mais une vente conjointe lente abaisse le total~; rejoindre une liquidation gérée si elle est contraignante et
couvre les positions qui se recoupent, et se couvrir dans l'intervalle.
\iqlookfor{la structure à somme nulle de la course.}
\end{solution}

\begin{solution}{iq:fm:case-studies-ii:4}
Seulement ce dont le trading a besoin ce jour-là, ramené chaque jour, sous une limite fixée par ce que la société pourrait perdre sans difficulté, puisque le solde constitue une
créance non garantie tant que la ségrégation n'est pas vérifiée.
\iqlookfor{le solde comme exposition de crédit.}
\end{solution}

\begin{solution}{iq:fm:case-studies-ii:5}
Que la bourse peut annuler les transactions selon ses règles, comme l'a fait le marché du nickel~; ne rien enregistrer tant que les transactions ne sont pas confirmées comme valides.
\iqlookfor{pouvoirs d'annulation.}
\end{solution}

\begin{solution}{iq:fm:case-studies-ii:6}
Avec un impact permanent linéaire, la baisse totale est fixée à la moitié du coefficient d'impact multipliée par la position~; l'ordre détermine qui la supporte, et le
dernier vendeur en supporte le plus. Des ventes conjointes plus lentes ou des transactions de bloc peuvent abaisser le total.
\iqlookfor{l'argument intégral.}
\end{solution}
